% ============================================================================
%  PHỤ LỤC
% ============================================================================
\clearpage
\appendix
\renewcommand{\thesection}{\Alph{section}}
\titleformat{\section}{\normalfont\Large\bfseries}{Phụ lục \thesection.}{0.6em}{}

\section{Hướng dẫn tái lập}
\label{app:repro}

Toàn bộ kết quả tái lập được bằng các lệnh sau trên một máy có Python 3.12 (GPU không bắt buộc
cho phần đánh giá baseline và phân tích; bắt buộc trên thực tế cho fine-tune).

\begin{lstlisting}[language=bash,escapeinside={(*}{*)}]
# 1. (*\lstcmt{Môi trường và dữ liệu}*)
uv venv .venv --python 3.12
uv pip install --python .venv/bin/python -r requirements.txt
python scripts/fetch_data.py                   # UIT-ViQuAD 2.0 -> data/raw/

# 2. (*\lstcmt{Kiểm thử (không cần}*) GPU)
python -m pytest -m "not slow"

# 3. (*\lstcmt{Phân tích không cần mô hình}*)
python scripts/analyze_segmentation.py         # (*\lstcmt{biên từ + trần}*) EM
python scripts/build_stress_v2.py              # (*\lstcmt{dựng + kiểm định}*) stress-test v2

# 4. Fine-tune (dev (*\lstcmt{tách từ train; validation không được nhìn thấy)}*)
python scripts/finetune.py --model FacebookAI/xlm-roberta-base --out models/xlmr \
    --epochs 3 --max-length 384 --doc-stride 128 --max-answer-len 64
python scripts/finetune.py --model vinai/phobert-base-v2 --out models/phobert \
    --word-segmented --epochs 3 --max-length 256 --doc-stride 96 --max-answer-len 64
bash scripts/queue_night1.sh && bash scripts/queue_night2.sh   # seed 13, lr 1e-5, 256/96

# 5. (*\lstcmt{Chấm (mỗi hệ thống một lần, toàn bộ validation), chọn tau trên dev, chẩn đoán}*)
python scripts/run_eval.py --models abstain baseline xlmr phobert
python scripts/run_eval.py --models abstain baseline xlmr xlmr_seed13 phobert \
    phobert_seed13 xlmr_256 --dataset stress2
bash scripts/queue_score_windows.sh
python scripts/calibrate_thresholds.py --runs xlmr xlmr_seed13 xlmr_256 phobert \
    phobert_seed13 phobert_lr2e5 phobert_stable
python scripts/diagnose.py

# 5b. Stress-test v2 (*\lstcmt{ở tau chọn trên dev (không chọn gì trên}*) stress2)
for m in xlmr xlmr_seed13 phobert phobert_seed13 xlmr_256; do
  python scripts/score_windows.py --checkpoint models/$m --name $m --split stress2
done
python scripts/score_stress2_tuned.py --check-tau0 \
    --runs xlmr xlmr_seed13 phobert phobert_seed13 xlmr_256

# 6. (*\lstcmt{Báo cáo: số liệu và hình sinh từ results/, rồi biên dịch}*)
python scripts/make_report_numbers.py && python scripts/make_figures.py
cd report && latexmk -xelatex main.tex
\end{lstlisting}

\section{Kiểm định khung mã ban đầu}
\label{app:legacy}

Trước khi chạy bất kỳ thí nghiệm nào, nhóm rà soát khung mã ban đầu của đồ án (nay ở
\code{legacy/}). Rà soát cho thấy khung mã \emph{chưa từng được chạy} (không có tệp kết quả
hay checkpoint nào) nhưng tài liệu đi kèm đã trình bày số liệu. Các phát hiện được liệt kê để
người đọc biết vì sao khung mã này không được dùng, và vì sao báo cáo có quy tắc ``không con số
nào không có tệp kết quả''.

\begin{table}[H]
\centering
\caption{Các lỗi phát hiện trong khung mã ban đầu}
\label{tab:legacy}
\renewcommand{\arraystretch}{1.18}
\small
\begin{tabularx}{\textwidth}{@{}L L L@{}}
\toprule
\textbf{Lỗi} & \textbf{Hệ quả nếu chạy} & \textbf{Xử lý} \\
\midrule
Ba bộ số liệu kết quả khác nhau (README, slide, hằng \code{EXPECTED\_RESULTS}), không bộ nào có tệp nguồn & Báo cáo số không đo & Xoá toàn bộ; mọi số sinh từ \code{results/} \\
\code{compute\_metrics} so sánh chuỗi \code{"span\_\{s\}\_\{e\}"} thay vì văn bản đáp án; F1 được ``ước lượng'' bằng \code{EM $\times$ 0,95} & EM/F1 vô nghĩa & Giải mã span thật từ offset (Mục~\ref{sec:segtok}) \\
Giải mã PhoBERT giữ dấu ``\_'' của bước tách từ & Đáp án đúng bị chấm EM~0 & Cắt đáp án từ văn bản gốc; kiểm thử hồi quy \\
Ghi đè \code{max\_position\_embeddings} của PhoBERT & Hỏng embedding vị trí khi nạp & Bỏ; \code{max\_length} = 256 \\
\code{include\_stress\_test=True} mặc định: 60\% bộ stress-test bị trộn vào tập \emph{huấn luyện} & Rò rỉ: chấm trên dữ liệu đã học & Stress-test chỉ dùng để đánh giá \\
Bộ sinh stress-test: \code{\_find\_answer\_start} trả về 0 khi không tìm thấy đáp án; \code{\_wrap\_into\_squad} gán lại ngữ cảnh theo vòng cho mẫu không có ngữ cảnh & Phần lớn câu có đáp án không nằm trong ngữ cảnh: bộ cũ không đo được thứ nó tuyên bố đo & Dựng lại bộ stress-test từ validation, có kiểm định tự động (Mục~\ref{sec:audit}) \\
\bottomrule
\end{tabularx}
\end{table}

\section{Cấu trúc mã nguồn}
\label{app:code}

\begin{lstlisting}[escapeinside={(*}{*)}]
src/mrc/
  data.py, normalize.py, metrics.py     # (*nạp SQuAD-2.0, chuẩn hoá,*) EM/F1      (CS116)
  windowing.py, features.py, training.py, transformer_qa.py                     (CS116)
  baseline_tfidf.py, evaluate.py, tagging.py                                   (CS116)
  tokenization.py         # (*chọn tokenizer theo loại mô hình*)                   (CS221)
  segmented_tokenizer.py  # PhoBERT: (*tách từ + BPE + offset về văn bản gốc*)     (CS221)
  segmentation.py         # (*tương thích biên từ, trần*) EM                       (CS221)
  stress_v2.py            # (*dựng, kiểm định, chấm theo cặp*) stress-test v2       (CS221)
  diagnosis.py            # (*phân loại lỗi, từ chối,*) McNemar, bootstrap         (CS221)
  baselines.py            # (*luôn từ chối*)                                       (CS221)
scripts/   finetune.py run_eval.py analyze_segmentation.py build_stress_v2.py
           score_windows.py calibrate_thresholds.py score_stress2_tuned.py
           diagnose.py make_report_numbers.py make_figures.py
tests/     (*\NTests{} kiểm thử (\NTestsFast{} không cần GPU)*)
results/   eval_*.json predictions_*.json training_curve_*.json diagnosis_*.json ...
report/    (*mã LaTeX của báo cáo;*) generated/ do make_report_numbers.py sinh
legacy/    khung (*mã ban đầu (không dùng; xem Phụ lục*) B)
\end{lstlisting}

\section{Bảng truy vết số liệu}
\label{app:trace}

\begin{table}[H]
\centering
\caption{Nguồn gốc các nhóm số liệu trong báo cáo}
\label{tab:trace}
\renewcommand{\arraystretch}{1.2}
\small
\begin{tabularx}{\textwidth}{@{}L L@{}}
\toprule
\textbf{Số liệu} & \textbf{Nguồn} \\
\midrule
Thống kê dữ liệu, kích thước train/dev & \code{data/raw/*.json} qua \code{compute\_stats}, \code{split\_by\_context} (seed 42) \\
EM, F1, tách nhóm, độ trễ & \code{results/eval\_*\_validation.json} \\
Điểm trên stress-test v2 theo nhóm, ở $\tau$ dev (Bảng~\ref{tab:stress}) & \code{results/eval\_*\_stress2\_tuned.json} \\
Độ nhất quán theo cặp E2c/E3b/E5 (Bảng~\ref{tab:stresspairs}) & \code{by\_subset.*.pairs} trong cùng tệp \\
Điểm từng cửa sổ trên stress-test v2 & \code{results/windows\_*\_stress2.json} \\
Kiểm định stress-test v2 & \code{results/stress\_v2\_audit.json} \\
Tương thích biên từ, trần EM & \code{results/segmentation\_validation.json} \\
Phân loại lỗi, từ chối, theo độ dài đáp án, so sánh cặp & \code{results/diagnosis\_validation.json} \\
Ví dụ bất đồng, mẫu lỗi định tính & \code{results/disagreements\_phobert\_xlmr.json}, \code{results/error\_sample\_*.json} \\
Đường cong huấn luyện, siêu tham số & \code{results/training\_curve\_*.json}; nhật ký \code{results/logs/train\_*.log} \\
Mọi macro số và bảng sinh tự động & \code{scripts/make\_report\_numbers.py} $\to$ \code{report/generated/} \\
Hình & \code{scripts/make\_figures.py} $\to$ \code{report/figures/} \\
\bottomrule
\end{tabularx}
\end{table}
